%------------------------------------------------------------------------------%
% THE SPECTRAL THEOREM FOR SELF-ADJOINT OPERATORS                              %
%------------------------------------------------------------------------------%
% File  : Spectral_Theorem.tex                                                 %
% Author: Klaus Hoeflinger, Beethovenstrasse 11, D-73117 Wangen                %
% Email : khoeflinger@t-online.de                                              %
%------------------------------------------------------------------------------%

%------------------------------------------------------------------------------%
%                       DOCUMENT CLASS and INCLUDE FILES                       %
%------------------------------------------------------------------------------%
\documentclass[11pt,twoside]{article}
\usepackage{geometry}
\geometry{paperheight=241mm,
          paperwidth=152mm,
          top=22mm,
          bottom=17mm,
          left=20mm,
          right=20mm,
          textwidth=112mm,
          textheight=202mm,
          headheight=10mm,
          headsep=3mm,
          footskip=9mm,
          includehead,
          includefoot,
          marginparsep=0mm,
          marginparwidth=0mm}

\renewcommand{\baselinestretch}{0.945}\normalsize

\AtBeginDocument{\setlength\abovedisplayskip{2mm}}
\AtBeginDocument{\setlength\belowdisplayskip{2mm}}

%------------------------------------------------------------------------------%
% define title formats                                                         %
%------------------------------------------------------------------------------%
\usepackage{titlesec}

\renewcommand{\thesection}{\arabic{section}}
\renewcommand{\sfdefault}{FiraSans}
\renewcommand{\ttdefault}{pcr}
\renewcommand{\rmdefault}{antiqua}

\titleformat{\part}[display]
{\normalfont \large \rmfamily \bfseries \centering}
{\small{\thepart.}}{2pt}{}

\titleformat{\section}[runin]
{\normalfont \rmfamily \bfseries}
{\thesection}{1pt}{.\;}[.]

%\titleformat{\section}
%{\normalfont \bfseries \small \filcenter}
%{\thesection.\;}{0mm}{}

\titleformat{\subsection}[runin]
{\normalfont \itshape}
{}{0mm}{}

\titleformat{\subsubsection}[runin]
{\normalfont \itshape}
{}{}{}{}

\titleformat{\paragraph}[runin]
{\normalfont}
{}{}{}{}

\titleformat{\subparagraph}[runin]
{\normalfont}
{}{}{}{}

\titlespacing*{\part}          { 0pt}{ 0pt}{ 8pt}
\titlespacing*{\section}       { 0pt}{ 7pt}{ 4pt}
\titlespacing*{\subsection}    { 0pt}{14pt}{ 6pt}
\titlespacing*{\subsubsection} { 0pt}{ 6pt}{12pt}
\titlespacing*{\paragraph}     { 0pt}{ 6pt}{12pt}
\titlespacing*{\subparagraph}  { 0pt}{ 0pt}{12pt}

%------------------------------------------------------------------------------%
% define enumerate parameters                                                  %
%------------------------------------------------------------------------------%
\usepackage{enumitem}
\setlength{\parindent}{3pt}
\setlength{\parskip}  {0pt}

%\setlist{noitemsep}

\setlist[enumerate]{
         leftmargin=12pt,
         rightmargin=0pt,
         itemsep=1pt,
         itemindent=3pt,
         topsep=3pt,
         labelsep=4pt,
         partopsep=0pt,
         labelwidth=0pt,
         listparindent=0pt,
         parsep=0pt}

\setlist[itemize]{
         leftmargin=12pt,
         rightmargin=0pt,
         itemsep=0pt,
         itemindent=1pt,
         topsep=1pt,
         labelsep=4pt,
         partopsep=1pt,
         labelwidth=0pt,
         listparindent=0pt,
         parsep=0pt}

%------------------------------------------------------------------------------%
% define packages                                                              %
%------------------------------------------------------------------------------%
\usepackage{upref}
\usepackage{amssymb}
\usepackage{slantsc}
\usepackage{amsmath}
\usepackage{amsthm}
\usepackage{thmtools}
\usepackage{amscd}
\usepackage{verbatim}
\usepackage{latexsym}
\usepackage{multicol}
\usepackage{graphicx}
\usepackage{setspace}
\usepackage[ngerman,english]{babel}
\usepackage[protrusion=true,expansion=true]{microtype}
\usepackage{tikz}
\usetikzlibrary{arrows,arrows.meta,decorations.markings}
\usepackage{mathrsfs}
\usepackage{amsfonts}
\usepackage{datetime2}
\usepackage{textgreek}
\usepackage{hyperref}
\usepackage{pifont}
\usepackage{relsize}
%\usepackage{biblatex}
\relsize{-0.05}
\usepackage[skip=0mm]{parskip}
\usetikzlibrary{decorations.pathmorphing}

\usepackage{mlmodern}
\usepackage[T5,T1]{fontenc}
\usepackage{ragged2e}

%\usepackage{mathabx}
%\usepackage{times}

%\usepackage{fontspec}
%\setmainfont{Nimbus Roman No9 L}
%\usepackage[T1]{fontenc}

%------------------------------------------------------------------------------%
% define page layout                                                           %
%------------------------------------------------------------------------------%
\usepackage{fancyhdr}
\pagestyle{fancy}
\setlength{\headheight}{30.0pt}
\renewcommand{\headrulewidth}{0pt}

\AtBeginDocument{
  \addtolength\abovedisplayskip{-0.0\baselineskip}
  \addtolength\belowdisplayskip{-0.0\baselineskip}
}

\fancyhead{}
\fancyfoot{}
\addtolength{\headwidth}{0.02\marginparwidth}

\makeatletter
\let\ps@plain\ps@fancy
\makeatother

\renewcommand\sectionmark[1]
{
 \def\sectionname{#1}
 \markright{\thesection #1}
}

\makeatletter
\@addtoreset{section}{part}
\makeatother

%------------------------------------------------------------------------------%
% define footnote without number                                               %
%------------------------------------------------------------------------------%
\newcommand{\myfootnote}[1]{%
\renewcommand{\thefootnote}{}%
\footnotetext{#1}%
\renewcommand{\thefootnote}{\arabic{footnote}}%
}

%------------------------------------------------------------------------------%
% define numbering in equations                                                %
%------------------------------------------------------------------------------%
\numberwithin{equation}{section}

%------------------------------------------------------------------------------%
% define newtheorem                                                            %
%------------------------------------------------------------------------------%
\declaretheoremstyle[
headfont=\scshape, 
headindent = 0mm, %\parindent,
%postheadhook = {\hspace{0mm}\newline},
%postheadspace = \newline, 
spaceabove = 3mm, 
spacebelow = 3mm,
numberwithin=section,
name=Theorem]{khMATH}
\declaretheorem[style=khMATH]{theorem}

\declaretheoremstyle[
headfont=\bfseries, 
headindent = 0mm, %\parindent,
%postheadhook = {\hspace{0mm}\newline},
%postheadspace = \newline, 
spaceabove = 3mm, 
spacebelow = 3mm,
numberwithin=part,
name=Theorem]{khMATH1}
\declaretheorem[style=khMATH1]{theorem1}

\declaretheoremstyle[
headfont=\scshape, 
headindent = 0mm, %\parindent,
%postheadhook = {\hspace{0mm}\newline},
%postheadspace = \newline, 
spaceabove = 3mm, 
spacebelow = 3mm,
numberwithin=section,
name=Corollary]{khMATH1}
\declaretheorem[style=khMATH1]{corollary}

%            name         counter     label              numeration-mode
\theoremstyle{plain}
\newtheorem*{introduction}            {\sc Introduction}
\newtheorem {standard}                {}                 [section]
\newtheorem*{definition}              {\sc Definition}
\newtheorem{satz1}                    {\sc Satz}
\newtheorem{lemma}                    {\sc Lemma}
\newtheorem*{proposition}             {\sc Proposition}
%\newtheorem{corollary}               {\sc Corollary}
\newtheorem*{corollar}                {\sc Korollar}
\newtheorem*{remark}                  {\sc Remark}
\newtheorem*{remarks}                 {\sc Remarks}
\newtheorem*{example}                 {Example}
\newtheorem*{examples}                {Examples}
\newtheorem*{theo}                    {\sc Theorem}

%------------------------------------------------------------------------------%
% define tabular                                                               %
%------------------------------------------------------------------------------%
\usepackage{tabularx}
\newcolumntype{L}[1]{>{\raggedright\arraybackslash}p{#1}}
\newcolumntype{C}[1]{>{\centering\arraybackslash}  p{#1}} 
\newcolumntype{R}[1]{>{\raggedleft\arraybackslash} p{#1}} 

%------------------------------------------------------------------------------%
% define \sh, \zB                                                              %
%------------------------------------------------------------------------------%
\def\dsh{{d.\,h.}}
\def\ie{{i.\,e.}}
\def\zB{z.\,B.}
\renewcommand{\abstractname}{ }

%------------------------------------------------------------------------------%
% change default spacing for cases                                             %
%------------------------------------------------------------------------------%
\makeatletter
\renewcommand*\env@cases[1][1.6]{%
  \let\@ifnextchar\new@ifnextchar
  \left\lbrace
  \def\arraystretch{#1}%
  \array{@{}l@{\quad}l@{}}%
}
\makeatother

\newenvironment{rcases}
  {\left.\begin{aligned}}
  {\end{aligned}\right\rbrace}

%------------------------------------------------------------------------------%
% change default for \predisplaypenalty                                        %
%------------------------------------------------------------------------------%
\predisplaypenalty=990
\allowdisplaybreaks \numberwithin{equation}{section}

%------------------------------------------------------------------------------%
% general notations                                                            %
%------------------------------------------------------------------------------%
\def\*{\star}
\def\={\,=\,}
\def\+{\,+\,}
\def\xsubset{{\,\subset\,}}
\def\xtimes{{\,\times\,}}
\def\xeof{{\,\in\,}}
\def\eq{{\,=\,}}
\def\xto{{\,\to\,}}
\def\eqdef{\,:=\,}
\def\:{{\,:\,}}
\def\ele{{\,\in\,}}
\def\exists{\mbox{ exists }}
\def\and{\mbox{ and }}
\def\for{\mbox{ for }}
\def\fa{\mbox{ for all }}
\def\fe{\mbox{ for each }}
\def\qfa{\quad \mbox{ for all }}
\def\df{\,\dotfill}
\def\iff{if and only if }
\def\wrt{with respect to}

\makeatletter
\newcommand \Df {\leavevmode \cleaders \hb@xt@ .70em{\hss .\hss }\hfill \kern \z@}
\makeatother

\def\sql{\mathfrak{a}}               % sesquilinear form a
\def\DF{B}                           % Dirichlet form a
%------------------------------------------------------------------------------%
% number sets and special elements                                             %
%------------------------------------------------------------------------------%
\def\P{\mathbb{H}}                   % half plane of $\C$
\def\J{I}                            % interval I
\def\N{{\mathbb{N}}}                 % unsigned integers
\def\Q{{\mathbb{Q}}}                 % rational integers
\def\Z{{\mathbb{Z}}}                 % integers
\def\R{{\mathbb{R}}}                 % real numbers
\def\Rn{{\mathbb{R}}^n}              % real numbers in Rn
\def\Rd{{\mathbb{R}}^d}              % real numbers in Rn
\def\C{{\mathbb{C}}}                 % complex numbers
\def\F{{\mathbb{K}}}                 % arbitrary field
\def\Torus{{\mathbb{T}}}             % complex circle line
\def\T{\tau}                         % upper bound of interval (0,t)
\def\sign#1{\operatorname{sign}(#1)} % sign of a number

%------------------------------------------------------------------------------%
% set theory                                                                   %
%------------------------------------------------------------------------------%
\def\G{G}                            % domain $G$
\def\clG{\overline{\G}}              % closure of $G$
\def\bndG{\partial \G}               % boundary of $G$
\def\setA{\mathcal{A}}               % set A
\def\setB{\mathcal{B}}               % set B
\def\setC{\mathcal{C}}               % set C
\def\setM{M}                         % set M
\def\setN{N}                         % set N
\def\setS{\mathcal{S}}               % set S
\def\famF{\mathfrak{F}}              % family of sets
\def\indI{\mathcal{I}}               % index set I
\def\pot#1{\mathfrak{P}(#1)}         % power set

\def\relR{\mathbf{R}}                % relation R
\def\mapf{f}                         % mapping f
\def\mapg{g}                         % mapping g

\def\set#1#2{\{ \, #1 \,: \, #2 \, \}}
\def\tensor#1#2{#1 \otimes #2}
\def\Tens{\oplus}

\def\cal#1{{\mathcal{#1}}}
\def\aut#1{\operatorname{Aut}(#1)}

\def\cross#1#2{#1 {\,\times\,} #2}
\def\multicross#1#2{#1 \times  \ldots \times  #2}
\def\multitens#1#2 {#1 \otimes \ldots \otimes #2}

%------------------------------------------------------------------------------%
% group theory                                                                 %
%------------------------------------------------------------------------------%
\def\1{{\mathsf{1}}}                % unit element of a monoid
\def\0{{\mathsf{0}}}                % unit element of an Abelian monoid
\def\kernel#1{\mathfrak{N}(#1)}     % kernel of a homomorphism
\def\cokernel#1{\mathfrak{C}{(#1)}} % cokernel of a homomorphism
\def\Hom#1{\operatorname{Hom}(#1)}

%------------------------------------------------------------------------------%
% linear algebra                                                               %
%------------------------------------------------------------------------------%
\def\vecV{\mathit{V}}               % vector space V
\def\vecH{\mathit{H}}               % vector space H
\def\vecW{\mathit{W}}               % vector space W
\def\vecU{\mathit{U}}               % subspace U
\def\convC{\mathcal{C}}             % convex set C
\def\basH{\mathfrak{B}}             % Hamel base B

\def\LO#1{\mathscr{L}(#1) }         % algebra of linear transformations

\def\scalar#1#2{ \langle #1,#2 \rangle }
\def\trace#1{\operatorname{tr}(#1)}
\def\dim#1{{\operatorname{dim} \, #1 }}
\def\codim#1{{\operatorname{codim} \, #1 }}
\def\dim{\operatorname{dim}}
\def\codim{\operatorname{codim}}
\def\ind{\operatorname{ind}}

\def\genG{\cal{G}}

\def\algA{\cal{A}}
\def\algB{\cal{B}}
\def\algU{\cal{U}}

\def\idI{\cal{I}}

%------------------------------------------------------------------------------%
% topology                                                                     %
%------------------------------------------------------------------------------%
\def\compK{{K}}          % compact set C
\def\openO{\mathcal{O}}  % open set O
\def\openU{\mathcal{U}}  % open set U
\def\U{\mathcal{U}}      % neighbourhood U
\def\V{\mathcal{V}}      % neighbourhood V
\def\d{\mathit{d}}       % metric function d

\def\interior#1{\mathring{#1}}
\def\closure#1{{\overline{#1}}}

\def\topT{\mathcal{T}}   % topology T
\def\topX{X}             % topological space X
\def\topY{Y}             % topological space Y
\def\topZ{Z}             % topological space Z
\def\metrS{M}            % metric space M

\def\grpG{G}             % topological group G
\def\grpA{A}             % topological group A
\def\grpB{B}             % topological group B
\def\grpC{C}             % topological group C
\def\grpH{H}             % topological group H
\def\grpU{U}             % topological group U
\def\grpV{V}             % topological group V
\def\topG{G}             % topological group G
\def\topH{H}             % topological group H
\def\topU{U}             % topological subgroup U

\def\subS{\cal{S}}
\def\riR{R}              % ring R

%------------------------------------------------------------------------------%
% calculus                                                                     %
%------------------------------------------------------------------------------%
\def\supp#1{\operatorname{supp}(#1)}
\def\singsupp#1{\operatorname{sing \, supp}(#1)}

\def\re{\operatorname{Re}}
\def\im{\operatorname{Im}}
\def\grad{\operatorname{grad}}

%\def\abs#1 {\left|#1\right|}
\def\abs#1 {\lvert #1 \rvert}
%\def\abs#1 {\left|\,#1\,\right|}

%------------------------------------------------------------------------------%
% measure theory                                                               %
%------------------------------------------------------------------------------%
\def\salgA{\Sigma}               % sigma-algebra S
\def\salgB{\cal{B}}              % sigma-algebra B
\def\Borel#1{\cal{B}_{#1}}       % Borel-algebra 
\def\LB{\lambda}                 % Lebesgue-Borel measure

%------------------------------------------------------------------------------%
% function spaces                                                              %
%------------------------------------------------------------------------------%
\def\Lp{L^p(\G)}                 % spaces L_p
\def\Lq{L^q(\G)}                 % spaces L_q
\def\Lh{L^2(\G)}                 % spaces L_2
\def\Li{L^{\infty}(\G)}          % spaces L_infty
\def\Lc{L_{loc}^1(\G)}           % spaces L_1^{loc}
\def\Ck{C^k(\G)}                 % spaces C_k

\def\BV#1{BV(#1) }               % set of functions of bounded variation
\def\AC#1{AC(#1) }               % set of absolutely continuous functions

\def\MP#1{\mathscr{M}^+(#1) }    % set of positive measures
\def\MS#1{\mathscr{M}^-(#1) }    % vector space of finite signed measures
\def\MC#1{\mathscr{M}(#1) }      % vector space of complex measures
\def\MCR#1{\mathscr{M}_{r}(#1) } % vector space of regular complex measures

\def\SobW{W^{k,p}(\G)}                % Sobolev space W
\def\SobO{\mathring{W}^{k,p}(\G)}     % Sobolev space W
%\def\WN#1#2{\mathring{W}^{#1}(#2) }   % Sobolev space W
\def\WN#1#2{W_0^{#1}(#2) }   % Sobolev space W
%\def\HN#1#2{\mathring{H}^{#1}(#2) }   % Sobolev space H
\def\HN#1#2{H_0^{#1}(#2) }   % Sobolev space H
%------------------------------------------------------------------------------%
% functional analysis                                                          %
%------------------------------------------------------------------------------%
\def\snorm#1{\lVert #1 \rVert}
\newcommand{\norm}[1]{\left\lVert #1 \right\rVert}

\def\topV{V}                       % topological vector space V
\def\topW{W}                       % topological vector space W
\def\normS{X}                      % normed space

\def\banX{\mathit{X}}              % Banach space X
\def\banY{\mathit{Y}}              % Banach space Y
\def\banZ{\mathit{Z}}              % Banach space Z
\def\dbanX{\mathit{X^*}}           % dual space of Banach space X
\def\dbanY{\mathit{Y^*}}           % dual space of Banach space Y
\def\dbanZ{\mathit{Z^*}}           % dual space of Banach space Z

\def\banA{\mathit{A}}              % Banach algebra A
\def\banB{\mathit{B}}              % Banach algebra B
\def\banC{\mathit{C}}              % C*-algebra C


\def\domain#1{\mathfrak{D}(#1)}    % domain
\def\range#1{\mathfrak{R}(#1)}     % range

\def\isoJ{\mathfrak{J}}            % isometry J
\def\repA{{\cal{A}_{\sql}}}        % representation operator A
\def\linS{{S}}                     % linear operator S
\def\linG{{G}}                     % linear operator G
\def\linL{{L}}                     % linear operator L
\def\linT{{T}}                     % linear operator A
\def\linA{\mathit{A}}              % linear operator A
\def\linB{{B}}                     % linear operator B
\def\linM{{M}}                     % linear operator M
\def\linU{{U}}                     % linear operator U
\def\linI{{I}}                     % linear operator I
\def\A{\cal{A}}                    % linear differential operator A
\def\BDO{\cal{B}}                  % linear boundary differential operator B
\def\linU{{U}}                     % unitary operator U
\def\linP{{P}}                     % projection P
\def\compA{k}                      % compact operator k
\def\linFR{f}                      % finite rank operator f
\def\fredA{A}                      % Fredholm operator F

\def\HAM{\cal{H}}                  % Hamilton operator
\def\PA{\partial^{\alpha}}         % elementary differential operator
\def\DO{\cal{A}}                   % linear differential operator
\def\BC{\cal{B}}                   % boundary operator
\def\SLO{\cal{L}_{p,q}}            % Sturm-Liouville operator
\def\MO{\cal{L}_{M}}               % momentum operator

\def\HS#1{S(#1)           }        % algebra of Hilbert-Schmidt operators
\def\FR#1{F(#1)           }        % algebra of finite rank operators
\def\TC#1{N(#1)           }        % algebra of trace class operators
\def\BU#1{\mathscr{U}(#1) }        % algebra of unitary linear operators
\def\BO#1{\mathscr{L}(#1) }        % algebra of bounded linear operators
\def\CO#1{\mathscr{C}(#1) }        % algebra of compact linear operators
\def\CL#1{\mathscr{C}(#1) }        % algebra of compact linear operators
\def\FO#1{\Phi(#1) }               % set of Fredholm operators
\def\FK#1#2{\Phi_{#2}(#1) }        % set of Fredholm operators with index k

\def\hilH{\mathit{H}}              % Hilbert space H
\def\clU{\mathit{U}}               % closed subspace U

\def\orthO{\mathfrak{O}}           % orthonormal system O
\def\orthS{\mathfrak{S}}           % orthonormal system S
\def\basB{\mathfrak{B}}            % basis B of a Hilbert space

\def\GL#1 {\mathbf{GL}(#1)}        % linear group
\def\OG#1 {\mathbf{O}(#1)}         % linear group
\def\SOG#1 {\mathbf{SO}(#1)}       % linear group

%------------------------------------------------------------------------------%
% distribution theory                                                          %
%------------------------------------------------------------------------------%
\def\disF{F}                  % distribution F
\def\disG{G}                  % distribution G
\def\disH{H}                  % distribution H
\def\disU{U}                  % distribution U
\def\disT{T}                  % distribution T
\def\SF{C^{\infty}(\G)}       % space of smooth functions
\def\TF{C_c^{\infty}(\G)}     % space of compactly supported smooth functions
\def\TFRd{C_c^{\infty}(\R^d)} % space of test functions
\def\SRd{\mathscr{S}(\R^d)}   % Schwartz space
\def\SR1{\mathscr{S}(\R)}     % Schwartz space
\def\B1#1{C^1(#1) }           % space of continuously differentiable functions
\def\Bm#1{C^m(#1) }           % space of continuously differentiable functions
\def\Ck{C^k(\G)}
\def\TD{\mathscr{S}'(\R^d)}   % space of temperated distributions
\def\E{\cal{E}}               % space of distributions with compact support
\def\D{\mathscr{D}}           % D
\def\FT{\mathcal{F}}          % Fourier transformation F
\def\FT{\mathscr{F}}          % Fourier transformation F

%------------------------------------------------------------------------------%
% other definitions                                                            %
%------------------------------------------------------------------------------%
\def\Proof{\textsc{Proof}}
\def\FRECHET{Fr\'{e}chet}
\def\ARZELA{Arzel\'{a}}
\def\GARDING{G\r{a}rding}
\def\HOELDER{H\"older}
\def\SCHROEDINGER{Schr\"odinger}
\def\JOERGENS{J\"orgens}
\def\WUEST{W\"uest}
\def\KAEHLER{K\"aehler}
\def\HOEFLINGER{H\"oflinger}
\def\POINCARE{Poincar\'{e}}
\def\FA{Functional Analysis}

\def\Author   {K.\,{\HOEFLINGER}}
\def\AuthorUC {K.\,HOEFLINGER}
\def\Version  {1.0}
\def\Email    {khoeflinger@t-online.de}


%\renewcommand\thesection{\Roman{section}}

\titleformat{\section}
{\normalfont \bfseries \filcenter}
{\thesection.\;}{0mm}{}

\titleformat{\subsection}%[runin]
{\normalfont \itshape}
{}{}{}

\titlespacing*{\section}   {0pt}{12pt}{4pt}
\titlespacing*{\subsection}{0pt}{6pt}{0pt}

%------------------------------------------------------------------------------%
%                                BEGIN OF DOCUMENT                             %
%------------------------------------------------------------------------------%
\begin{document}
\thispagestyle{empty}

\def\Title{The Spectral Theorem for Self-Adjoint Operators}

\centerline{{\textit{Topics in Spectral and Semigroup Theory}}}
\vspace{22pt}
\centerline{\large{\textbf{\Title}}}
\vspace{4pt}
\centerline{\small{{\textsc{\Author}}}}
\vspace{15pt}

\fancyhead[C] {\small{\textsc{\Title}}}
\fancyhead[OL]{\small{\thepage}}
\fancyhead[ER]{\small{\thepage}}

\myfootnote
{
  Version {\Version} from \today;
  Email:  \textit{khoeflinger@\,t-online.de}
}

\begin{abstract}
This exposition is a review on the spectral theorem for not necessarily bounded
self-adjoint linear operators in Hilbert space and its proof.
Some applications of the spectral theorem are given.
\end{abstract}

%------------------------------------------------------------------------------%
\section{Introduction}                                                         %
%------------------------------------------------------------------------------%
The spectral theorem for self-adjoint linear operators in his most general form
was formulated and proved in the twentieth of the last century
primarily by the mathematicians \textsc{John von\,Neumann}
and \textsc{Harvey\;M. Stone}
and is one of the oldest and most important results in the upcoming mathematical
theory of \textit{functional analysis}.

\subparagraph{}
Nowadays there are several formulations of this theorem available,
and they are in fact equivalent to each another. 
Let $\hilH$ be a complex Hilbert space.
Then the classical form of the theorem quotes
that each self-adjoint linear operator $\linA$ acting in $\hilH$
is a \textit{spectral operator},
that is,
there is a representation of the form
\begin{equation}\label{ST_1}
\scalar{\linA \, x}{y} \=
\int_{\,\R} \lambda \, d\scalar{P_{\lambda} \,x}{y}
\quad \mbox{for } x \xeof \domain{\linA} \mbox{ and } y \xeof \hilH,
\end{equation}
and the \textit{domain of definition} $\domain{\linA}$
of the operator $\linA$ is given by
\begin{equation}\label{ST_2}
\domain{\linA} \=
\{\, x \xeof \hilH \:
    \int_{\,\R}
\lambda^2 \, d \scalar{P_{\lambda} \,x}{x} < \infty\,\}.
\end{equation}

\subparagraph{}
Here by $P_{\lambda}$ is meant a member
of any family of \textit{projections} in $\hilH$
parametrized by the real numbers
and usually denoted $\{P_{\lambda}\}_{\lambda \in \R}$.
The integral in \eqref{ST_1} means to integrate
the identity mapping $\lambda \xeof \R \mapsto \lambda$
{\wrt} the family $\{P_{\lambda}\}_{\lambda \in \R}$ in the sense of vector-valued
Riemann-Stieltjes integration and Bochner integration,
respectively.

\subparagraph{}
We basically shall go ahead in this text step by step
by giving well-understandable descriptions of all the upcoming concepts.

%------------------------------------------------------------------------------%
\section{Basic Facts on Operators in Hilbert Space}                            %
%------------------------------------------------------------------------------%
Let $\hilH$ be a complex Hilbert space,
where we do not suppose that it is a separable one.
Given elements $x,y \xeof \hilH$,
let $\scalar{x}{y} \xeof \C$ be the inner product
and $\norm{x} := \sqrt{\scalar{x}{x}}$ be the derived norm of $x \xeof \hilH$.

\subparagraph{}
We are dealing with \textit{linear} operators
$\linA \:\domain{\linA} \xsubset \hilH \xto \hilH$,
so they are fulfilling the computation rules
\[
\linA(\lambda x+ \mu y) = \lambda \linA x + \mu \linA y
\quad \quad (x,y \xeof \domain{\linA}, \; \lambda,\mu \xeof \C).
\]
The set $\domain{\linA}$,
called the \textit{domain of definition} of $\linA$,
is assumed to be a linear subspace of $\hilH$.
It is also supposed that this set lies dense in $\hilH$,
so we may denote $\linA$ a \textit{densely defined} linear operator in $\hilH$.

\subparagraph{}
Since $\linA$ is densely defined,
its \textit{adjoint operator} $\linA^*$ is uniquely defined
due to \textsc{Riesz}'s representation theorem,
and the domain of definition $\domain{\linA^*}$
of $\linA^*$ consists of those elements $y \xeof \hilH$,
for which there is $y^* \xeof \hilH$ fulfilling
\[
\scalar{\linA \,x}{y} \eq \scalar{x}{y^*} \fa x \xeof \domain{\linA}.
\]
We then define $\linA^* y := y^*$ and denote the linear operator $\linA^*$
the \textit{Hilbert adjoint} of $\linA$.
It is well-known
that if $\linA$ is \textit{closable},
then $\linA^*$ is densely defined,
too,
hence the bi-adjoint operator $\linA^{**} := (\linA^*)^*$ is well-defined
and coincides with the \textit{closure} $\overline{\linA}$ of $\linA$.
Adjoint operators are basically closed,
which in fact is not hard to show.

\paragraph{}
A linear operator $\linA$ acting in $\hilH$
and with domain of definition $\domain{\linA}$ dense in $\hilH$
is called \textit{symmetric},
if the adjoint $\linA^*$ forms an extension of $\linA$
(notation: $\linA \prec \linA^*$),
which is equivalent to
\[
\scalar{\linA x}{y} \= \scalar{x}{\linA y} \qfa x,y \in \domain{\linA}.
\]

\subparagraph{}
Furthermore,
a linear operator $\linA$ is called \textit{self-adjoint},
if $\linA \eq \linA^*$.
Clearly,
each self-adjoint operator is symmetric.
For a bounded linear operator,
where $\domain{\linA} \eq \hilH$ can be assumed,
there is obviously no difference between symmetry and self-adjointness.

\paragraph{}
We finally want to introduce some basic concepts from \textit{spectral theory}.
Let again $\linA$ be a linear operator in $\hilH$,
but not necessarily densely defined.
The a complex number $\zeta \xeof \C$ is called a \textit{regular value}
{\wrt} $\linA$,
if $\linA-\zeta \1$
(from now on shortly denoted $\linA-\zeta$)
is bijective and the \textit{inverse operator}
\[
(\linA-\zeta)^{-1} \:\hilH \to \domain{\linA} \subset \hilH
\]
is bounded.
The \textit{bounded inverse theorem} implies that
if $\linA-\zeta$ is closed and bijective,
then $(\linA-\zeta)^{-1}$ proves to be bounded,
so boundedness of the inverse does not have to be postulated in addition
to ensure regularity of $\zeta$.

\subparagraph{}
The set $\rho(\linA)$ consisting of all regular values of $\linA$ is called
the \textit{resolvent set} of $\linA$,
whereas the complementary set $\sigma(\linA) := \C \setminus \rho(\linA)$
is called the \textit{spectrum} of $\linA$.
Moreover,
the mapping
\[
R(\linA,\cdot): \rho(\linA) \to \BO{\hilH}\footnote{
\,$\BO{\hilH}$ is the Banach algebra of bounded linear operators
$\linA\:\hilH \xto \hilH$ with operator norm
$\norm{\linA} := \sup_{\norm{x}=1} \norm{\linA x}$
for $\linA \xeof \BO{\hilH}$.}
\]
is called the \textit{resolvent function}
or briefly the \textit{resolvent} of $\linA$.
It is an operator-valued function, which is even holomorphic
(in the sense of vector-valued functions).

\paragraph{}
The spectrum $\sigma(\linA)$ and the resolvent $R(\linA,\cdot)$
of a self-adjoint linear operator $\linA$ enjoy nice properties,
which we collect to the following

\begin{theorem}
\,
\begin{itemize}[leftmargin=22pt]
\item[(a)]
$\linA$ is self-adjoint \iff $\sigma(\linA) \xsubset \R$.
\item[(b)]
If $\linA$ is self-adjoint and $\zeta \xeof \rho(\linA)$,
then $\im \zeta \neq 0$ by (a) and
\begin{equation}\label{RESOLVENT_ESTIMATE}
\norm{R(\linA,\zeta)} \leq \frac{1}{\abs{\im \zeta} }.
\end{equation}
\end{itemize}
\end{theorem}

%------------------------------------------------------------------------------%
\section{Projection Operators}                                                 %
%------------------------------------------------------------------------------%
Let $\setA \xsubset \hilH$ be a \textit{closed} linear subspace.
Then the \textit{projection theorem} claims that
for each $x \xeof \hilH$ there is exactly one $x_0 \xeof \setA$
such that
\[
\norm{ x - x_0 } \= \inf \, \{ \, \norm {x-y} \: y \in \setA \,\}.
\]
The element $x_0$ is called the \textit{projection} of $x$ onto $\setA$
and minimizes the distance between $\setA$ and $x$.
The associated mapping
$\pi_{\setA} \: x \xeof \hilH \mapsto x_0 \xeof \setA$
is linear,
denoted the \textit{projection} of $\hilH$ onto $\setA$.
Alternatively,
a \textit{projection} is a linear operator
$\pi \:\hilH \xto \hilH$ fulfilling the conditions $\pi  \circ \pi  \eq \pi $
and $\pi \eq \pi ^*$.
Projections are important examples of symmetric rsp. self-adjoint operators.

\subparagraph{}
By the projection theorem,
projections are in one-to-one relationship to closed linear subspaces in $\hilH$.
Hence it is nearby
to call two projections $\pi_1,\pi_2$ \textit{orthogonal} to each other,
if their related closed subspaces have this property,
and in this case the identity
$\pi_1 \circ \pi_2 \eq \pi_2 \circ \pi_1 \eq \0$ holds,
where $\0$ is the projection related
with the trivial subspace $\{\0\} \subset \hilH$.

\paragraph{} %- properties of projections
In general,
a projection mapping $\pi$ has the following properties:

\begin{itemize}[leftmargin=12pt,itemsep=5pt]
\item[1.]
$\pi$ is bounded with $\norm{\pi} \leq 1$,
hence $\pi \xeof \BO{\hilH}$.\\
If $\setA_{\pi} := \pi(\hilH) \eq \{\0\}$,
then $\norm{\pi} \eq 0$,
else $\norm{\pi}=1$.
The only possible eigenvalues of $\pi$ are the real numbers $0$ and $1$.

\item[2.]
Let $\setA_1,\setA_2$ be closed subspaces of $\hilH$
with $\setA_1 \subseteq \setA_2$,
then for the associated projections the equation
$\pi_{\setA_1} \circ \pi_{\setA_2} \eq
 \pi_{\setA_2} \circ \pi_{\setA_1} \eq \pi_{\setA_1}$
holds.
In this case,
the operator $\pi_{\setA_2} - \pi_{\setA_1}$ is also a projection,
and related subspace is given by the set
$\setA_2 \ominus \setA_1 := \setA_2 \cap \setA_1^{\bot}$.

\item[3.]
$\pi \eq \pi^{*}$,
that is,
$\pi$ is symmetric.
Conversely,
each symmetric and idempotent linear operator is a projection.

\item[4.]
If $\setA$ is a \textit{finite}-dimensional subspace of $\hilH$,
then $\setA$ is closed,
and the projection $\pi_{\setA}$ can be represented
by means of a set of base vectors $e_1,\ldots,e_n$ in $\setA$, namely
\[
\pi_{\setA} x \= \sum_{i=1}^n \scalar{x}{e_i} \, e_i.
\]

\item[5.]
Let $\mathscr{P}(\hilH)$ be the set of projections in $\hilH$,
then there is a partial order relation $\prec$ on $\mathscr{P}(\hilH)$,
defined by means of inclusion between the related closed subspaces:
\[
\pi_{\setA_1} \prec \pi_{\setA_2} \quad
: \Longleftrightarrow \quad \setA_2 \subseteq \setA_1.
\]
\end{itemize}

%------------------------------------------------------------------------------%
\section{Tools from Measure Theory and Complex Analysis}                       %
%------------------------------------------------------------------------------%
The proof of the spectral theorem given in the next section requires some
additional utilites from other mathematical areas.
We especially need the the concept of \textit{complex measure} and the conept
of a \textit{Herglotz function} from complex analysis.

\paragraph{} %-- complex measures
In contrast to positive measures,
a \textit{complex measure} on the measurable space $(X,\salgB)$
is a $\sigma$-additive set function $\mu \: \salgB \xto \C$.
The condition $\mu(\emptyset) \eq 0$ is fulfilled by itself
and need not be part of the definition.

\subparagraph{}
When dealing with complex measures,
it is possible to define the concept of \textit{variation}.
More precisely,
the variation $\abs{\mu} $ of a complex measure $\mu$
is that mapping $\abs{\mu} \: \salgB \xto [0,\infty]$,
which assigns to each measurable set $A$ the value
\[
\abs{\mu} (A) \eqdef
\sup \, \sum_{n=1}^{\infty} \, \norm{\mu(A_n)},
\]
where the sequences $(A_n)$ may be any mostly countable disjoint
decompositions of $A$ in measurable subsets.
It is a surprising result
that for each complex measure $\mu$
the variation $\abs{\mu} $ is actually a \textit{finite} positive measure.

\subparagraph{}
The \textit{total variation} of a complex measure $\mu$ is defined
to be the variation of the entire set $X$,
that is,
$\abs{\mu} := \abs{\mu} (X)$.
Since $\abs{\mu} < \infty$,
each complex measure is said to be of \textit{bounded variation}.

\paragraph{} %-- Herglotz functions
What we are using from complex analysis, is a property
of so-called \textit{Herglotz functions}.

\subparagraph{}
Let $\H^+ := \{z \xeof \C \: \im z > 0\}$ be the upper half plane of $\C$.
Then a function $\Phi \: \H^+ \xto \C$ is called a {Herglotz-function}
(sometimes also called a \textit{Nevanlinna function},
if $\Phi$ is holomorphic and $\im \Phi(z) \geq 0$.
Every such function admits an integral representation of the form
\begin{equation}\label{Nevanlinna}
\Phi(z) \= C + D\,z +
\int_{\R} (\frac{1}{\lambda-z} -\frac{\lambda}{1+\abs{\lambda} ^2}) \, d\mu(\lambda)
\end{equation}
with $C \xeof \R$, $D \geq 0$ and $\mu$ being a Borel measure satisfying the estimate
\[
\int_{\R} \frac{1}{1+\abs{\lambda} ^2} \,d\mu < \infty.
\]
Reversely,
every function defined by \eqref{Nevanlinna} is Herglotz.
The Borel measure $\mu$ can be recovered from $\Phi$ by
\begin{equation}\label{INVERSION_FORMULA}
\mu((\lambda_1,\lambda_2]) \eq \lim_{\delta \to 0} \, \lim_{\epsilon \to 0}
\int_{\lambda_1+\delta}^{\lambda_2+\delta}
\im \Phi(\lambda + i \epsilon) \, d\lambda,
\end{equation}
called the \textit{Stieltjes inversion formula} for Herglotz functions
and their integral representations.

\subparagraph{}
These considerations are meaningful,
since the function
\[
\Phi_A(\zeta,x) \eqdef \scalar{R(\linA,\zeta)x}{x}
\]
is Herglotz for every self-adjoint linear operator $\linA$
and every $x \xeof \domain{\linA}$.
Moreover, due to \eqref{RESOLVENT_ESTIMATE} we have
\[
\abs{\Phi_A(\zeta,x)} \leq \frac{\norm{x}^2}{\abs{\im \zeta} }
\quad \quad (x \xeof \domain{\linA}, \; \zeta \xeof \H^+)\,,
\]
and with respect to the integral representation \eqref{Nevanlinna}
we can simplify equation \eqref{Nevanlinna} to
\begin{equation}\label{Nevanlinna_2}
\Phi_A(\zeta,x) \=
\int_{\R} \frac{1}{\lambda-\zeta} \, d\mu_{A,x}(\lambda),
\end{equation}
where of course the Borel measure $\mu_{A,x}$ depends on the operator $\linA$
and the fixed chosen vector $x \xeof \domain{\linA}$.

%------------------------------------------------------------------------------%
\section{Spectral Measures and Spectral Operators}                             %
%------------------------------------------------------------------------------%
Let $\Borel{\R}$ be the Borel $\sigma$-algebra of the real line.
Then a \textit{projection-valued measure}
is a mapping $P \: \Borel{\R} \xto \BO{\hilH}$,
which fulfills the following properties:

\begin{itemize}[leftmargin=12mm]
\item[(P-1)]
$P(\Omega)$ is a projection operator for each $\Omega \xeof \Borel{\R}$;

\item[(P-2)]
$P(\R) \eq \1_{\hilH}$, the identity operator in $\hilH$;

\item[(P-3)]
$P$ is \textit{$\sigma$-additive},
that is,
if $\Omega_1, \Omega_2,\ldots$ are measurable subsets of the real line
and $\Omega := \bigsqcup_{\,n \in \N} \Omega_n$,
then
\[
P(\Omega) \= \lim_{n \xto \infty} \, \sum_{k=1}^n \, P(\Omega_k),
\]
where this kind of convergence has to be understood
in the sense of the strong operator topology on the algebra $\BO{\hilH}$,
that is,
$P(\Omega) x \eq \lim_{n \xto \infty} \, \sum_{k=1}^n \, P(\Omega_k) x$
for each $x \xeof \hilH$.
\end{itemize}

\subparagraph{}
The projection-valued measure $P$ is also called
a \textit{spectral measure} and owns the following properties:

\begin{itemize}[leftmargin=12pt]
\item[1.]
$P(\emptyset) \eq 0$.
\item[2.]
$P(\R \setminus \Omega) \eq \1_{\hilH} - P(\Omega)$.
\item[3.]
$P(\Omega_1 \sqcup \Omega_2) \eq P(\Omega_1) + P(\Omega_2)$
\item[4.]
$P(\Omega_1 \cap \Omega_2) \eq P(\Omega_1) \circ P(\Omega_2)$
\item[5.]
$\Omega_1 \subset \Omega_2  \Longrightarrow P(\Omega_1) \prec P(\Omega_2)$
\end{itemize}

\subparagraph{}
Notice that
intervals of the form $(-\infty,\lambda]$ are Borel-measurable,
so for each $\lambda \xeof \R$ 
the projection $\linP((-\infty,\lambda])$ is well-defined.
Let $\cal{P}(\lambda) := P((-\infty,\lambda))$.
The one-parameter family $\{\cal{P}(\lambda)\}_{\lambda \in \R}$
is called a \textit{resolution of the identity}
and uniquely determines the original projection-valued measure $\linP$.

\paragraph{} %-- sprectral integrals
In the sequel
it is our aim to introduce a \textit{spectral integral}
for measurable functions $f\:\R \xto \C$
according to a projection-valued measure $P$
and a resolution of the identity $\cal{P}$,
respectively.

\subparagraph{}
Let $B(\R)$ be the linear space of \textit{bounded} measurable functions
$f\:\R \xto \C$ and endowed with the supremum norm
\[
\norm{f}_{\infty} \eqdef \sup_{x \in \R} \, \abs{f(x) } < \infty,
\]
which in fact proves to be a Banach space.
Any function $e \xeof B(\R)$ is called \textit{simple},
if it adopts only finitely many values $e_1,\ldots, e_N$.
The norm of $e \xeof S(\R)$ is given by
$\norm{e}_{\infty} := \max \, \{\abs{e_1} ,\ldots, \abs{e_N} \}$.
The set
\[
S(\R) := \{\,f\:\R \xto \C: f \mbox{ is simple}\,\}
\]
constitutes a dense linear subspace of $B(\R)$,

\paragraph{}
After having these preliminaries available,
the setup of the spectral integral will be carried out
by the following five steps:

\begin{itemize}[leftmargin=12pt,itemsep=2mm]
\item[1.]
First of all,
each spectral measure $P$ gives rise
to a complex and real-valued Borel measure on the real line.
For arbitrary $x,y \xeof \hilH$
let the complex measure $\mu_{x,y}$ be defined by
\[
\begin{cases}[1.2]
\quad \mu_{x,y}\: \Borel{\R} \xto \C,\\
\quad \Omega \mapsto \mu_{x,y}(\Omega) := \scalar{x}{P(\Omega)y}.\\
\end{cases}
\]
Due to the symmetry of the projection $P(\Omega)$,
the derived measure
$\mu_x(\Omega) := \mu_{x,x}(\Omega)$
for $\Omega \xeof \Borel{\R}$ is real-valued.

\item[2.]
We next make sense of the expression $\int_{\R} f \, dP$
for measurable functions $f\:\R \xto \C$.
Let us start with the case, where $f$ is \textit{simple},
that is,
$f$ takes values in the finite set $\{f_1,\ldots,f_N\}$.
Define
\[
\int_{\R} f \, d\linP \eqdef \sum_{n=1}^N f_n \linP(\chi^{-1}(f_n)),
\]
where $\chi^{-1}(f_n)$ is the preimage of $f_n$ in $\Borel{\R}$.
We further mention
that
\[
I_{0,\linP}\:f \xto \int_{\R} f \, d\linP
\]
is a continuous linear mapping from the normed space
$S(\R)$ consisting of simple functions
to the Banach space $\BO{\hilH}$ with operator norm $\norm{I_{0,P}} \eq 1$.

\item[3.]
We shall continue with the integration
of \textit{bounded measurable functions} {\wrt} $P$.
Due to the B.L.T. theorem\footnote{
\,The B.L.T. theorem claims that each continuous and densely defined linear operator
may be extended to a continuous linear operator on the whole of the domain space,
where the operator norm is preserved.},
there exists a unique extension $I_P$ from $I_{0,P}$
to the closure $B(\R) \eq \overline{S(\R)}$,
where $\norm{I_P} \eq \norm{I_{0,P}}$.
The mapping $I_P$ is also linear and fulfills
\[
\begin{cases}[1.2]
\quad I_P(\1_{\R}) = \1_{\hilH}\\
\quad I_P(fg) = I_P(f) \circ I_P(g)\\
\quad I_P(\overline{f}) = I_P(f)^*\\
\end{cases}
\quad \quad f,g \in B(\R),
\]
where $\1_{\R}$ is the identity in $\R$,
thus $I_P$ even forms a C*-algebra homomorphism between
the function algebra $B(\R)$ and the operator algebra $\BO{\hilH}$.

\item[4.]
We will also define $I_P$
for an \textit{arbitrary measurable function} $f\:\R \xto \C$.
But this is only possible for those elements $x \xeof \hilH$ fulfilling
\begin{equation}\label{PVM_OPERATOR}
\int_{\R} \abs{f} ^2 \, d \mu_x < \infty.
\end{equation}
The set $D_f$ of elements $x \xeof \hilH$ satisfying \eqref{PVM_OPERATOR}
forms a linear subspace lying dense in $\hilH$.
For $f$ we introduce an approximating sequence $\{f_n\}_{n \in \N}$
of functions $f_n \in B(\R)$ given by
\[
f_n \eqdef \chi_{ \{t \in \R: \abs{f(t)} \leq n\}} \cdot f
\quad (n \in \N).
\]
One can show that $f_n$ is Cauchy in $L^2(\R)$,
hence $I_P(f_n)$ is Cauchy in $\hilH$,
where $I_P$ is the previously defined operator acting on the space $B(\R)$.

\item[5.]
To the end,
for a fixed measurable function $f$
we may define an operator $S(P,f)$ by virtue of
\begin{equation}\label{SPECTRAL_OPERATOR}
\int_{\R} f \, dP \: D_f \mapsto \hilH\,,
\quad \quad
x \in D_f \mapsto (\int_{\R} f \, d\linP\,) \, x \eqdef y\,,
\end{equation}
which enjoys the property
\[
(\int_{\R} f \, d\linP)^* \= \int_{\R} \overline{f} \, dP.
\]
The mapping $S(P,f)$ is called a \textit{spectral operator}
with data $P$ and $f$.
Especially for $f := \1_{\R}$ 
we obtain the self-adjoint spectral operator $\linS_P := \linS(P,\1)$
according to
\[
\domain{\linS_P} := D_{\1_{\R}},
\quad \linS_P x \eqdef (\int_{\R} \1_{\R} \, dP) \, x.
\]

\subparagraph{}
Let $\linS_P^*$ be the adjoint of $\linS_P$.
Then we have
\[
\linS_P^* = (\int_{\R} \1_{\R} \, dP)^* = \int_{\R} \overline{\1_{\R}} \, dP = \linS_P,
\]
which shows that $\linS_P$ is self-adjoint.
Rather similar,
it can be shown that operator $S(P,f)$ is self-adjoint iff $f$ is real-valued.
\end{itemize}

\paragraph{}
We denote a linear operator $\linA$ is \textit{spectrally decomposible},
if there is a spectral operator $\linS_P$ such that $\linA \eq \linS_P$.
Let $\linA$ be spectrally decomposible,
then for any measurable function $f\:\R \to \C$ we may define the operator
\[
f(\linA) := \int_{\R} f\,dP,
\]
where the domain of definition $\domain{f(\linA)}$ is given by
\[
\domain{f(\linA)} \=
\{ x \in \hilH \: \int_{\R} \abs{f} ^2\,d\mu_x(\lambda) < \infty \}.
\]

\subparagraph{}
For example,
let $f(\lambda) := e^{-\lambda}$,
then the operator $\exp(A) := \linS(P,f)$ proves to be self-adjoint.
If $f(\lambda) := e^{i \lambda}$, however,
then $\exp(iA)$ constitutes an unitary operator in $\hilH$.
Further operators like $\sqrt{\linA}$, $\sin(\linA)$, $\cos(\linA)$,
$\abs{\linA} $ and so on are defined in the same way.
It should be emphasized that the spectral operators
$\exp(iA)$, $\sin(\linA)$, $\cos(\linA)$ are everywhere defined in $\hilH$.

\subparagraph{}
To the end,
let $\linA$ be a spectrally decomposible operator and $\zeta \xeof \rho(\linA)$.
Then the inverse operator $R(\linA,\zeta) := (\linA-\zeta)^{-1}$ exists and we
obtain with $r_{\zeta}(\lambda) := (\lambda-\zeta)^{-1}$
\[
R(\linA,\zeta)
\= r_{\zeta}(A)
\= \int_{\R} r_{\zeta} \,dP
\= \int_{\R} \frac{1}{\lambda-\zeta} \, dP(\lambda)
\]
and for every $x \in \hilH$
\[
\Phi(\zeta,x) \eqdef 
\scalar{x}{R(\linA,\zeta)x} \=
\int_{\R} \frac{1}{\lambda-\zeta} \,d\mu_x(\lambda).
\]
The mapping $\Phi(\cdot,x) \: \H^+ \xto \H^+$ is a Herglotz function,
that is,
$\im \zeta \geq 0$ implies $\im \Phi(\zeta,x) \geq 0$.
This can be shown by
\[
\im \Phi(\zeta,x) = \im \int_{\R} \frac{1}{\lambda-\zeta} \,d\mu_x(\lambda) = 
\int_{\R} \frac{\im \zeta}{\abs{\lambda-\zeta} ^2} \,d\mu_x(\lambda) =
\im \zeta \cdot \phi(\zeta,x),
\]
where 
\[
\phi(\zeta,x) \eqdef \int_{\R} \frac{1}{\abs{\lambda-\zeta} ^2} \,d\mu_x(\lambda)
\, \geq \, 0.
\]

%------------------------------------------------------------------------------%
\section{Quotation and Proof of the Spectral Theorem}                          %
%------------------------------------------------------------------------------%
As already mentioned in the introductory section of this text,
the spectral theorem for unbounded self-adjoint operators
may be quoted as follows:

\begin{theorem}\label{SPECTRAL_THEOREM}
If a linear operator $\linA$ acting in a complex Hilbert space
and with domain of definition $\domain{\linA}$ is self-adjoint,
then there is projection-valued measure $P$ such that

\begin{equation}\label{ST_1}
\linA \=
\int_{\,\R} \lambda \, d\,P_{\lambda},
\end{equation}

where the linear subspace $\domain{\linA}$ is given by
\begin{equation}
\domain{\linA} \=
\{\, x \xeof \hilH \:
    \int_{\,\R}
\lambda^2 \, d \scalar{P_{\lambda} \,x}{x} < \infty\,\}.
\end{equation}
\end{theorem}

\subparagraph{}
In other words,
the theorem quotes that each self-adjoint operator $\linA$
is spectrally decomposible.
Finally,
since spectral operators of type $\linS(P,\1)$ are self-adjoint,
we may conclude that there is a one-to-one relationship between
the set of self-adjoint operators in $\hilH$ and
the set of projection-valued measures on $\hilH$.

\paragraph{}
The remaining part of this section deals with the proof
of our main theorem \ref{SPECTRAL_THEOREM}.
The basic idea is to construct a spectral measure $P$ for each Herglotz function
\begin{equation}\label{REP_1}
\Phi(\zeta,x)
\eqdef \scalar{R(\linA,\zeta)x}{x}
\= \int_{\R} \frac{1}{\lambda-\zeta} \, d\scalar{Px}{x}\,(d\lambda)
\end{equation}
for $x \in \domain{\linA}$ and $\im \zeta > 0$.
The function $\Phi(\cdot,x)$ is well-defined in $\H^+$,
since the spectrum of $\linA$ is a subset of the real line.
It can be shown that \eqref{REP_1} is equivalent to the "vector-valued"
variant
\begin{equation}\label{REP_2}
R(\linA,\zeta)x \= \int_{\R} \frac{1}{\lambda-\zeta} \, dP_{\lambda}x
\quad \quad (\im \zeta \neq 0).
\end{equation}
Once this representation is established,
we may obtain the spectral decomposition of $\linA$ as follows:
let $\linA'(\zeta)$ be the spectral operator defined by
\[
\linA'(\zeta) x \eqdef \int_{\R} \frac{1}{\lambda-\zeta} \, dP_{\lambda}x.
\]
Then $\linA'(\zeta)$ is the resolvent of $\linA$ by \eqref{REP_2},
but also the resolvent of the spectral operator 
\[
\linA' x \eqdef \int_{\R} \, dP_{\lambda}x,
\]
which implies $\linA \eq \linA'$, since the identity of resolvents causes
the identity of the operators themselves,
that is, \eqref{ST_1} must be valid.

\paragraph{}
So we have to show that the representation formula \eqref{REP_1}
holds for some spectral measure $P$.
Remember the integral representation of a Herglotz function
and the Stieltjes inversion formula \eqref{INVERSION_FORMULA},
quoting that there is a Borel measure $\mu$ such that
\begin{equation}\label{INVERSION_FORMULA_2}
\mu((\lambda_1,\lambda_2]) \eq \lim_{\delta \to 0} \, \lim_{\epsilon \to 0}
\int_{\lambda_1+\delta}^{\lambda_2+\delta}
\im \scalar{R(\linA,\lambda + i \epsilon)x}{x} \, d\lambda.
\end{equation}

We first want to calculate the term
\[
\Psi(x,\delta,\lambda_1,\lambda_2) \eqdef
\lim_{\epsilon \to 0} \int_{\lambda_1+\delta}^{\lambda_2+\delta}
\im \scalar{R(\linA,\lambda + i \epsilon)x}{x} \, d\lambda
\]
in order to compute
\[
\mu(\lambda_1,\lambda_2]) =
\lim_{\delta \to 0} \Psi(x,\delta,\lambda_1,\lambda_2)
\]
in a second step.
In fact, we have
\[
\Psi(x,\delta,\lambda_1,\lambda_2) 
\=
\lim_{\epsilon \to 0} \int_{\lambda_1+\delta}^{\lambda_2+\delta}
\im \Phi(\lambda+ i \epsilon,x) \, d\lambda
\]
\[
\=
\lim_{\epsilon \to 0} \int_{\lambda_1+\delta}^{\lambda_2+\delta}
\int_{\R} \frac{\epsilon}{\zeta-(\lambda +i \epsilon)} \, \mu_x \, d\zeta \, d\lambda,
\]
which by \textsc{Fubini}'s integration formula may be written as
\[
\lim_{\epsilon \to 0} \int_{\R}
\int_{\lambda_1}^{\lambda_2} \frac{\epsilon}{(\zeta-\lambda)^2 + \epsilon^2} \, d\lambda \, \mu_x \, d\zeta.
\]

\subparagraph{}
The calculation of the inner integral yields
\[
\int_{\lambda_1}^{\lambda_2} \frac{\epsilon}{(\zeta-\lambda)^2 + \epsilon^2} \, d\lambda \=
\arctan(\frac{\zeta-\lambda}{\epsilon}) \bigg|_{\,\lambda_1}^{\,\lambda_2},
\]
hence
\[
\lim_{\epsilon \to 0} 
\int_{\lambda_1}^{\lambda_2} \frac{\epsilon}{(\zeta-\lambda)^2 + \epsilon^2} \, d\lambda \=
\pi \cdot \chi_{[\lambda_1,\lambda_2]}
\]
and
\[
\Psi(x,\delta,\lambda_1,\lambda_2) \=
\pi \cdot \chi_{[\lambda_1,\lambda_2]} \, \mu_x \, d\zeta.
\]
In conclusion,
it follows by the dominated convergence theorem
\[
\mu((\lambda_1,\lambda_2]) \= \pi \cdot
\lim_{\delta \to 0} \int_{\lambda_1+\delta}^{\lambda_2+\delta} \chi_{[\lambda_1,\lambda_2]} \, \mu_x \, d\zeta.
\= 
\pi \cdot \int_{\lambda_1}^{\lambda_2} \mu_x \, d\zeta.
\]

%------------------------------------------------------------------------------%
\section{Other Quotations and Proofs}                                          %
%------------------------------------------------------------------------------%

%------------------------------------------------------------------------------%
\section{Examples and Applications}                                            %
%------------------------------------------------------------------------------%

%------------------------------------------------------------------------------%
%                                 BIBLIOGRAPHY                                 %
%------------------------------------------------------------------------------%
\newpage
\thispagestyle{empty}

\vspace{8mm}
\begin{thebibliography}{99}
\begin{small}

\bibitem{RS1} M.\,Reed, B.\,Simon:
\textit{Modern Methods of Mathematical Physics I: Functional Analysis}.
Academic Press Inc., 1980.

\bibitem{RN} F.\,Riesz, B.\,Sz.-Nagy:
\textit{{Functional Analysis}}.
Springer, 1956.

\bibitem{Yo}  K.\,Yosida:
\textit{Functional Analysis}.
Springer, 1980.

\end{small}
\end{thebibliography}

%------------------------------------------------------------------------------%
%                               END OF DOCUMENT                                %
%------------------------------------------------------------------------------%
\end{document}
